\section{A conservative uniform reassembly certificate}
\label{sec:uniform-reassembly}

The coefficient-specific ratio is the sharpest statement available
from the triangle inequality.  For an archive it is often convenient
to compare every term with the same global raw normalization.

\begin{theorem}[Raw-normalized reassembly]
\label{thm:raw-reassembly}
Let \(D_{\rm raw}>0\), \(\Gamma_{\rm pre}>0\),
\(0\le\eta<1\), and \(\chi\ge0\).  Suppose a literal coefficient
\(z\) satisfies
\begin{align}
 \|\cA z\|^2
 &\ge(1-\eta)\Gamma_{\rm pre}D_{\rm raw},
 \label{eq:represented-lower}\\
 \langle D_{\cM}z,z\rangle
 &\le\chi\Gamma_{\rm pre}D_{\rm raw}.
 \label{eq:missing-diagonal-upper}
\end{align}
Then
\begin{equation}
 \|\cB z\|^2
 \ge
 \delta_{\rm sharp}\Gamma_{\rm pre}D_{\rm raw},
 \qquad
 \delta_{\rm sharp}
 :=
 \left(\sqrt{1-\eta}
       -\sqrt{\widehat\kappa\chi}\right)_+^2.
 \label{eq:sharp-reassembly}
\end{equation}
In particular,
\begin{equation}
 \boxed{
 \|\cB z\|^2
 \ge
 \delta_{\rm reasm}\Gamma_{\rm pre}D_{\rm raw},\qquad
 \delta_{\rm reasm}
 :=
 \left[
 1-2\sqrt{\widehat\kappa\chi}
 -\widehat\kappa\chi-\eta
 \right]_+.}
 \label{eq:uniform-reassembly}
\end{equation}
\end{theorem}

\begin{proof}
\Cref{lem:missing-energy} and
\eqref{eq:missing-diagonal-upper} imply
\[
 \|\cM z\|
 \le
 \sqrt{\widehat\kappa\chi\,
       \Gamma_{\rm pre}D_{\rm raw}}.
\]
Together with \eqref{eq:represented-lower}, the reverse-triangle
inequality gives \eqref{eq:sharp-reassembly}.

It remains to compare the displayed constants.  If the bracket in
\eqref{eq:uniform-reassembly} is nonpositive, the claim follows from
nonnegativity.  Otherwise
\(\sqrt{1-\eta}>\sqrt{\widehat\kappa\chi}\), and expansion gives
\begin{align*}
 \delta_{\rm sharp}
 &=1-\eta
 -2\sqrt{(1-\eta)\widehat\kappa\chi}
 +\widehat\kappa\chi\\
 &\ge
 1-\eta-2\sqrt{\widehat\kappa\chi}
 -\widehat\kappa\chi
 =\delta_{\rm reasm}.
\end{align*}
\end{proof}

\begin{remark}[Why retain the weaker constant]
The sharper constant in \eqref{eq:sharp-reassembly} should be used
when all quantities are exact.  The conservative form
\eqref{eq:uniform-reassembly} separates the represented defect
\(\eta\), the cross-term cost \(2\sqrt{\widehat\kappa\chi}\), and an
explicit missing self-energy reserve
\(\widehat\kappa\chi\).  It is therefore convenient for a route
ledger, but it is not advertised as sharp.
\end{remark}

\begin{proposition}[Certificate stop condition]
\label{prop:reassembly-stop}
The data \(\rho(P)<\infty\) and
\(\Gamma_{\rm pre}>0\) alone do not imply
\(\delta_{\rm reasm}>0\).  A positive conclusion requires a
quantitative relation among \(\widehat\kappa\), \(\chi\), and \(\eta\).
If
\[
 1-2\sqrt{\widehat\kappa\chi}
 -\widehat\kappa\chi-\eta\le0,
\]
then \eqref{eq:uniform-reassembly} supplies no positive survival
estimate and this certificate branch must stop.  This failure is not
a proof that the actual reassembled energy vanishes.
\end{proposition}

\begin{proof}
The formula itself gives the second assertion.  For the first, take
\(\cM=-\cA\) on the one-dimensional span of \(z\).  Then
\(\cB z=0\), although both maps have finite Gram matrices and
\(\Gamma_{\rm pre}\) may be positive.
\end{proof}

The one-dimensional example also shows why a claim of ``the missing
part is constructible'' is weaker than an energy statement.  Exact
constructibility can produce exact destructive interference.
